\chapter{Introduction}
\label{ch:introduction}

Multilingual large language models (LLMs) such as Llama~3.1 \citep{dubey2024llama3}, Qwen~3 \citep{yang2025qwen3}, and Mistral \citep{jiang2023mistral} can perform tasks in languages that received little or no explicit supervision during instruction tuning \citep{brown2020gpt3}. However, this zero-shot cross-lingual transfer remains inconsistent across language families. While models frequently succeed in high-resource languages such as French or German, performance often degrades on low-resource or typologically distant languages such as Amharic or Burmese. The mechanisms governing when and why cross-lingual transfer succeeds remain an active area of study.

One prominent account is the \emph{pivot hypothesis}: during pre-training on multilingual text, transformer models form a shared, language-agnostic representation space (an \emph{interlingua}) in their intermediate layers, mapping semantically equivalent sentences close to one another regardless of surface language \citep{wendler2024llamas}. Under this hypothesis, internal representational alignment should correlate with downstream multilingual capability without requiring task-specific evaluation. Validating this claim requires probing representation spaces directly across diverse languages and network depths.

MEXA (Multilingual Evaluation via Cross-Lingual Alignment) \citep{kargaran2025mexa} provides a framework for this type of probing. Given parallel sentence pairs in a target and a pivot language, MEXA extracts hidden representations across layers, computes cosine similarity matrices, and applies a bidirectional precision-at-1 retrieval criterion designed to account for embedding anisotropy \citep{ethayarajh2019contextualized, hammerl2023anisotropy}. The resulting metric produces layer-wise alignment scores per language without requiring text generation, task-specific fine-tuning, or reference translations.

This thesis reproduces the original MEXA methodology and extends the evaluation across four dimensions: recent decoder families across parameter scales, sparse mixture-of-experts architectures, masked and sentence encoders, and non-English pivot languages. The objective is to characterize where and how cross-lingual alignment forms across modern model architectures.


% ============================================================
\section{Motivation}
\label{sec:intro-motivation}

\paragraph{Disparities in multilingual coverage.}
While LLMs support translation, multilingual search, and document processing, standard evaluation benchmarks cover only a small subset of the world's languages and writing systems \citep{bandarkar2024belebele, lai2023okapi}. Languages written in non-Latin scripts or with limited digital corpora face several obstacles: subword tokenizers heavily fragment their words into short pieces (the \emph{fertility tax}), their intermediate representations remain distant from the high-resource representation space, and their downstream accuracy drops substantially \citep{nllbteam2022}. Isolating the source of these disparities requires internal diagnostic tools that evaluate representations rather than surface task outputs alone.

\paragraph{Reference-free evaluation.}
Standard multilingual benchmarks conflate representation quality with prompt sensitivity, reasoning capacity, and generation dynamics. For example, a model might solve a translated reading comprehension question through surface memorization rather than aligned conceptual representations. Conversely, a model with well-aligned hidden states might fail on a benchmark because of decoding errors. Probing hidden-state geometry directly through reference-free metrics isolates representation quality from task-specific generation artifacts, enabling cleaner comparisons across architectures, training data mixes, and parameter scales.

\paragraph{Architectural shifts.}
Following the original MEXA study, newer model families have introduced different design choices, including Qwen~3 and~3.5 \citep{yang2025qwen3}, Mixtral \citep{jiang2024mixtral}, and Gemma~4. These families vary in pre-training data composition, vocabulary size, and routing mechanisms such as sparse mixture-of-experts (MoE). How these architectural choices influence the formation and quality of multilingual representations remains an open question that this thesis investigates.


% ============================================================
\section{Research Questions}
\label{sec:intro-research-questions}

This thesis addresses four core research questions:

\begin{description}
    \item[RQ1 (Scale and Training Recipe)] \emph{How does cross-lingual alignment vary with model size and pre-training data composition?} \\
    We evaluate decoder families across parameter sweeps (Qwen~3 from 0.6B to 8B, Gemma~4 from 2B to 31B) and compare different families at matched parameter scale (Llama~3.1~8B, Qwen~3~8B, and Apertus~8B) to isolate the effects of parameter count, vocabulary design, and training data mixtures.

    \item[RQ2 (Sparse Architectures)] \emph{Does sparse mixture-of-experts routing disrupt or preserve the shared interlingua?} \\
    We compare dense baselines against corresponding MoE architectures (Mistral~7B against Mixtral~8x7B and~8x22B; dense Gemma~4 against Gemma~4~26B-A4B) to evaluate whether routing specialization fragments language-agnostic representations across experts.

    \item[RQ3 (Pivot Independence)] \emph{Is the emergent interlingua language-agnostic or biased toward English?} \\
    We evaluate alignment using alternative pivot languages (French, German, Arabic, Chinese, and Basque) representing different script types, resource levels, and language families. We then apply mean-centering to evaluate whether removing language-specific centroid offsets stabilizes alignment across pivots.

    \item[RQ4 (Predictive Validity)] \emph{Does internal representational alignment correlate with downstream multilingual task performance?} \\
    We compare layer-wise MEXA scores against performance on the Belebele reading comprehension benchmark \citep{bandarkar2024belebele} and the m-ARC reasoning benchmark \citep{lai2023okapi} across models and individual languages to assess whether internal alignment serves as an effective proxy for functional multilingual capabilities.
\end{description}


% ============================================================
\section{Contributions}
\label{sec:intro-contributions}

The primary contributions of this thesis are:

\begin{enumerate}
    \item \textbf{Reproduction and baseline verification.} We reproduce the original MEXA evaluation on Llama~3.1~8B and Mistral~7B~v0.3 under identical configurations, verifying pipeline correctness and establishing a baseline for subsequent comparisons.

    \item \textbf{Broad architectural coverage.} We evaluate MEXA across more than 25 models spanning four architecture classes: causal decoders (Qwen~3, Qwen~3.5, Apertus, Gemma~4), sparse MoE models (Mixtral~8x7B and~8x22B, Gemma~4~26B-A4B), masked language models (XLM-RoBERTa, Glot500 \citep{imani2023glot500}), and contrastive sentence encoders (LaBSE \citep{feng2022labse}, mE5 \citep{wang2024multilingual}).

    \item \textbf{Non-English pivot and geometric analysis.} We demonstrate that cross-lingual alignment is not solely an artifact of an English pivot. Evaluating five typologically diverse pivots shows consistent alignment profiles, and centering representation centroids yields peak alignment scores of $\mu_{\text{Max}} \ge 0.89$ across all tested pivots.

    \item \textbf{Fertility tax and script analysis.} We evaluate tokenizer fertility across all models and datasets, finding substantial negative correlations ($r \in [-0.46, -0.62]$) between subword fragmentation and alignment scores, with script family acting as a dominant factor in per-language performance.

    \item \textbf{Interactive analysis interface.} We implement an interactive React dashboard for inspecting layer trajectories, fertility correlations, retrieval metrics, and embedding projections across models, languages, and evaluation setups.
\end{enumerate}


% ============================================================
\section{Thesis Outline}
\label{sec:intro-thesis-outline}

The remainder of this thesis is structured as follows:

\begin{description}
    \item[Chapter~\ref{ch:background}: Background] covers multilingual transformer architectures, tokenization methods, cross-lingual transfer mechanisms, the role of pivot languages, evaluation strategies (task-based vs.\ reference-free), embedding space anisotropy, and the parallel evaluation corpora (FLORES-200 and the Parallel Bible Corpus).

    \item[Chapter~\ref{ch:methodology}: Methodology] details the four stages of the MEXA pipeline (hidden state extraction, cosine similarity calculation, bidirectional retrieval scoring, and layer aggregation) alongside extensions such as representation mean-centering and alternate pivot selection. It also details the 25+ evaluated models, evaluation setups, multi-level analysis protocols, and compute infrastructure.

    \item[Chapter~\ref{ch:results}: Results and Discussion] presents empirical findings, covering the baseline reproduction, decoder scaling trends, MoE routing behavior, non-English pivot experiments, per-language fertility analyses, and downstream correlations with Belebele and m-ARC.

    \item[Chapter~\ref{ch:conclusion}: Conclusion] summarizes the findings, examines methodological limitations, and outlines directions for future research.
\end{description}
